%%=============================================================================
%% Bijlage: studiemateriaal voor de gebruikersstudie
%%
%% Dit bestand voegt de standalone PDF's van het PDF-pack en de SUS-vragenlijst
%% als bijlagen in de bachelorproef toe via de pdfpages-package.
%%
%% Gebruik:
%%   1. Compileer eerst de standalone PDF's:
%%        xelatex pdfpack-blokA.tex
%%        xelatex pdfpack-blokB.tex
%%        xelatex sus-vragenlijst.tex
%%   2. Activeer de %%=============================================================================
%% Bijlage: studiemateriaal voor de gebruikersstudie
%%
%% Dit bestand voegt de standalone PDF's van het PDF-pack en de SUS-vragenlijst
%% als bijlagen in de bachelorproef toe via de pdfpages-package.
%%
%% Gebruik:
%%   1. Compileer eerst de standalone PDF's:
%%        xelatex pdfpack-blokA.tex
%%        xelatex pdfpack-blokB.tex
%%        xelatex sus-vragenlijst.tex
%%   2. Activeer de %%=============================================================================
%% Bijlage: studiemateriaal voor de gebruikersstudie
%%
%% Dit bestand voegt de standalone PDF's van het PDF-pack en de SUS-vragenlijst
%% als bijlagen in de bachelorproef toe via de pdfpages-package.
%%
%% Gebruik:
%%   1. Compileer eerst de standalone PDF's:
%%        xelatex pdfpack-blokA.tex
%%        xelatex pdfpack-blokB.tex
%%        xelatex sus-vragenlijst.tex
%%   2. Activeer de %%=============================================================================
%% Bijlage: studiemateriaal voor de gebruikersstudie
%%
%% Dit bestand voegt de standalone PDF's van het PDF-pack en de SUS-vragenlijst
%% als bijlagen in de bachelorproef toe via de pdfpages-package.
%%
%% Gebruik:
%%   1. Compileer eerst de standalone PDF's:
%%        xelatex pdfpack-blokA.tex
%%        xelatex pdfpack-blokB.tex
%%        xelatex sus-vragenlijst.tex
%%   2. Activeer de \input{bijlagen/bijlage-materialen} regel in
%%      DeVuystSilasBP.tex.
%%
%% De pdfpages-package wordt hier geladen omdat de hogentreport-class hem zelf
%% niet meelaadt; dat heeft geen invloed op andere onderdelen.
%%=============================================================================

\usepackage{pdfpages}

\chapter{Studiemateriaal voor de gebruikersstudie}
\label{ch:bijlage-materialen}

Tijdens de evaluatie (hoofdstuk~\ref{ch:evaluatie}) krijgt elke deelnemer hetzelfde leerstof aangeboden in twee condities: de Proof-of-Concept (PoC) en een traditioneel \emph{PDF-pack}. Na afloop wordt de bruikbaarheid van de PoC gemeten met een Nederlandstalige System Usability Scale.
Deze bijlage bundelt de exact gebruikte documenten zodat de studie reproduceerbaar blijft.

\section*{Bijlage B.1: PDF-pack Conceptblok A (levels 1--3)}
\addcontentsline{toc}{section}{Bijlage B.1: PDF-pack Conceptblok A}
\includepdf[pages=-,scale=0.95,pagecommand={}]{bijlagen/pdfpack-blokA.pdf}

\section*{Bijlage B.2: PDF-pack Conceptblok B (levels 4--6)}
\addcontentsline{toc}{section}{Bijlage B.2: PDF-pack Conceptblok B}
\includepdf[pages=-,scale=0.95,pagecommand={}]{bijlagen/pdfpack-blokB.pdf}

\section*{Bijlage B.3: SUS-vragenlijst (NL)}
\addcontentsline{toc}{section}{Bijlage B.3: SUS-vragenlijst}
\includepdf[pages=-,scale=0.95,pagecommand={}]{bijlagen/sus-vragenlijst.pdf}
 regel in
%%      DeVuystSilasBP.tex.
%%
%% De pdfpages-package wordt hier geladen omdat de hogentreport-class hem zelf
%% niet meelaadt; dat heeft geen invloed op andere onderdelen.
%%=============================================================================

\usepackage{pdfpages}

\chapter{Studiemateriaal voor de gebruikersstudie}
\label{ch:bijlage-materialen}

Tijdens de evaluatie (hoofdstuk~\ref{ch:evaluatie}) krijgt elke deelnemer hetzelfde leerstof aangeboden in twee condities: de Proof-of-Concept (PoC) en een traditioneel \emph{PDF-pack}. Na afloop wordt de bruikbaarheid van de PoC gemeten met een Nederlandstalige System Usability Scale.
Deze bijlage bundelt de exact gebruikte documenten zodat de studie reproduceerbaar blijft.

\section*{Bijlage B.1: PDF-pack Conceptblok A (levels 1--3)}
\addcontentsline{toc}{section}{Bijlage B.1: PDF-pack Conceptblok A}
\includepdf[pages=-,scale=0.95,pagecommand={}]{bijlagen/pdfpack-blokA.pdf}

\section*{Bijlage B.2: PDF-pack Conceptblok B (levels 4--6)}
\addcontentsline{toc}{section}{Bijlage B.2: PDF-pack Conceptblok B}
\includepdf[pages=-,scale=0.95,pagecommand={}]{bijlagen/pdfpack-blokB.pdf}

\section*{Bijlage B.3: SUS-vragenlijst (NL)}
\addcontentsline{toc}{section}{Bijlage B.3: SUS-vragenlijst}
\includepdf[pages=-,scale=0.95,pagecommand={}]{bijlagen/sus-vragenlijst.pdf}
 regel in
%%      DeVuystSilasBP.tex.
%%
%% De pdfpages-package wordt hier geladen omdat de hogentreport-class hem zelf
%% niet meelaadt; dat heeft geen invloed op andere onderdelen.
%%=============================================================================

\usepackage{pdfpages}

\chapter{Studiemateriaal voor de gebruikersstudie}
\label{ch:bijlage-materialen}

Tijdens de evaluatie (hoofdstuk~\ref{ch:evaluatie}) krijgt elke deelnemer hetzelfde leerstof aangeboden in twee condities: de Proof-of-Concept (PoC) en een traditioneel \emph{PDF-pack}. Na afloop wordt de bruikbaarheid van de PoC gemeten met een Nederlandstalige System Usability Scale.
Deze bijlage bundelt de exact gebruikte documenten zodat de studie reproduceerbaar blijft.

\section*{Bijlage B.1: PDF-pack Conceptblok A (levels 1--3)}
\addcontentsline{toc}{section}{Bijlage B.1: PDF-pack Conceptblok A}
\includepdf[pages=-,scale=0.95,pagecommand={}]{bijlagen/pdfpack-blokA.pdf}

\section*{Bijlage B.2: PDF-pack Conceptblok B (levels 4--6)}
\addcontentsline{toc}{section}{Bijlage B.2: PDF-pack Conceptblok B}
\includepdf[pages=-,scale=0.95,pagecommand={}]{bijlagen/pdfpack-blokB.pdf}

\section*{Bijlage B.3: SUS-vragenlijst (NL)}
\addcontentsline{toc}{section}{Bijlage B.3: SUS-vragenlijst}
\includepdf[pages=-,scale=0.95,pagecommand={}]{bijlagen/sus-vragenlijst.pdf}
 regel in
%%      DeVuystSilasBP.tex.
%%
%% De pdfpages-package wordt hier geladen omdat de hogentreport-class hem zelf
%% niet meelaadt; dat heeft geen invloed op andere onderdelen.
%%=============================================================================

\usepackage{pdfpages}

\chapter{Studiemateriaal voor de gebruikersstudie}
\label{ch:bijlage-materialen}

Tijdens de evaluatie (hoofdstuk~\ref{ch:evaluatie}) krijgt elke deelnemer hetzelfde leerstof aangeboden in twee condities: de Proof-of-Concept (PoC) en een traditioneel \emph{PDF-pack}. Na afloop wordt de bruikbaarheid van de PoC gemeten met een Nederlandstalige System Usability Scale.
Deze bijlage bundelt de exact gebruikte documenten zodat de studie reproduceerbaar blijft.

\section*{Bijlage B.1: PDF-pack Conceptblok A (levels 1--3)}
\addcontentsline{toc}{section}{Bijlage B.1: PDF-pack Conceptblok A}
\includepdf[pages=-,scale=0.95,pagecommand={}]{bijlagen/pdfpack-blokA.pdf}

\section*{Bijlage B.2: PDF-pack Conceptblok B (levels 4--6)}
\addcontentsline{toc}{section}{Bijlage B.2: PDF-pack Conceptblok B}
\includepdf[pages=-,scale=0.95,pagecommand={}]{bijlagen/pdfpack-blokB.pdf}

\section*{Bijlage B.3: SUS-vragenlijst (NL)}
\addcontentsline{toc}{section}{Bijlage B.3: SUS-vragenlijst}
\includepdf[pages=-,scale=0.95,pagecommand={}]{bijlagen/sus-vragenlijst.pdf}
